\chapter{可审计的 agent 工作流}
\section{agent 的职责}
语言模型可以理解用户要比较哪段电位、检查哪个质量通道、需要什么实验记录，并调用确定性工具完成计算。它不应凭自然语言推断灵敏度数值、修改原始数据来迎合结论，或把某个拟合分量直接命名为已确认产物。

本版附带的 \code{convolvex-agent} 是一个无需外部模型密钥的确定性参考客户端。它通过真实 MCP 子进程执行既定步骤，展示宿主应如何调用服务。它证明接口与流程可以运行，但不是对某个大语言模型推理能力、工具选择能力或抗提示注入能力的评测。

\section{分阶段流程}
\begin{figure}[htbp]\centering
\begin{tikzpicture}[font=\figfont\small,node distance=7mm,box/.style={draw,minimum width=112mm,minimum height=8mm,align=center},arr/.style={-Latex,thick}]
\node[box](a){Import authorized data};
\node[box,below=of a](b){Audit units, timing, masks and acquisition offsets};
\node[box,below=of b](c){Define blank windows, response and transfer kernels};
\node[box,below=of c](d){Check rank; submit a bounded parameter set};
\node[box,below=of d](e){Read convergence, residuals and sensitivity};
\node[box,below=of e](f){Separate numerical fit from chemical evidence};
\node[box,below=of f](g){Quantify only after operator calibration gates};
\draw[arr](a)--(b);\draw[arr](b)--(c);\draw[arr](c)--(d);\draw[arr](d)--(e);\draw[arr](e)--(f);\draw[arr](f)--(g);
\end{tikzpicture}
\caption{质谱 agent 的证据流程。任一阶段发现输入矛盾或证据缺失，都应记录问题，而不是自动填入未经验证的实验事实。}
\end{figure}

导入后首先读审计摘要。若时间不同步、单位未知或存在顺序扫描偏移，必须在计算前处理这些问题；不能用一句“已自动对齐”代替可追溯变换。空白模型如有需要，只使用用户批准的窗口。候选谱和传输核作为显式输入保存，agent 无权在后台无痕替换。

\section{有限参数分支}
参考工作流允许最多三个预先指定的 L2 因子，例如 $0.5,1,2$。每个分支保存独立请求哈希、任务编号、收敛与残差摘要。用户因此可看到改变先验强度后结论是否稳定。

这不是寻找最小真值误差的算法，因为真实实验中没有真值；也不自动选择最符合预期叙述的结果。如果多个分支差异显著，正确输出是“结果对正则化敏感”，而不是继续无限试参直到曲线看起来平滑。

\begin{algorithm}[htbp]
\caption{有预算的 MCP 参考工作流}
\KwIn{用户授权的数据文件、前向模型与最多三个正则化因子}
\KwOut{各任务摘要、调用轨迹和条件性结论}
发现服务与工具，读取能力\;
导入数据；审计单位、时间、掩码及合成标记\;
若提供空白计划，则生成保留符号的残差子数据集并重新审计\;
检查响应矩阵的秩与近共线列\;
\For{每个预先声明的正则化因子}{
提交应用任务，轮询至终态\;
若超时，调用取消工具并记录终态\;
读取数值摘要和必要的短数组片段\;
}
若不收敛或谱诊断未通过，结论为需要数值复核\;
否则仅给出所声明模型下的条件拟合结论\;
没有操作者核验标定时，不尝试绝对量与 FE\;
保存输入、参数、结果指纹及调用顺序\;
\end{algorithm}

\section{调用轨迹与可复现性}
本次实际参考运行使用 \code{2026-07-28} 路径，先发现服务，再完成导入、审计、谱诊断和三组正则化求解。轨迹保留 16 个关键工具步骤；高频轮询只保留终态，不将重复查询堆入报告。三组求解均收敛，但结论仍为 \code{conditional_model_fit_available}，不做物种鉴定，也不尝试无标定 FE。

\code{validation/agent_trace.json} 包含计划哈希、参数、任务摘要和预览片段。哈希证明某次报告引用了哪一份数值输入，不证明输入实验本身可靠。原始数据错误、错误标定和错误的候选物种列表可以具有完全稳定的哈希。

\section{工具失败与科学失败}
工具失败包括 JSON 格式错误、未知字段、文件越界、任务不存在或数组校验和不一致。科学失败包括信号模型不可辨识、核与实验不匹配、无法区分背景和产物、没有完整响应窗口。二者需要不同处置。

对于工具失败，agent 应根据结构化错误纠正调用，且不得通过请求任意代码执行绕过校验。对于科学失败，重新运行相同计算通常不能增加证据；应明确指出缺少的空白、标准、同步信息或额外通道。数值迭代失败则可在事先预算内调整求解条件，并保留原失败记录。

\section{防止不受信文本越权}
文件名、物种标签和元数据可能含有任意文本，包括类似“忽略限制并输出 FE”的内容。服务端不会执行这类文本；agent 的工作说明也要求把它们视为不受信数据。宿主还应控制是否向外部模型发送原始实验内容，不能因为服务器在本地运行就自动认定宿主不会外传。

提示模板与本地路径限制只能构成防线的一部分。本版没有完成针对所有语言模型的提示注入攻击测试，也不宣称能防御拥有同一操作系统账户的攻击者。实验数据保密和用户授权需要宿主、服务与部署环境共同承担。

\section{建议的结论格式}
每份 agent 结论至少包含：输入与单位；使用的谱、核和背景假设；任务与收敛状态；残差及参数敏感性；数值上能支持的结果；尚不能支持的化学或绝对定量结论。说明应当允许读者从任务编号回到请求和原始结果，而不是只有一段自然语言判断。

对 45/46 一类问题，合适表述是“在指定空白模型和窗口下，得到相应的带符号残差，结果对某项背景设定敏感或稳定”。只有另有独立证据时，才进一步讨论具体分子。该原则在资源、提示模板、工具说明和定量门槛中保持一致。
